\problemname{Egoökning}
\noindent

Efter flera år av intensivt programmeringsslipande har Tommy insett att han inte alls känner sig självsäker.

För att lösa detta vill Tommy öka sitt självförtroende och ego!

Initialt har Tommy $E$ ego.

Det finns $N$ ego-ökningsproblem (förkortat EP), där lösandet av ett EP ökar Tommys ego med 1.
Varje EP har ett ID som är ett unikt heltal från $1$ till $N$, och en svårighetsgrad $a_i$.

För att garantera att Tommy kan lösa ett EP måste hans ego vara minst lika stort som svårighetsgraden av EP:t.
Med andra ord måste $E \geq a_i$ gälla för att Tommy ska kunna lösa det $i$:te EP:t.

Tommy kan börja med att lösa vilket som helst av de givna $N$ EP:n.
Men efter att ha löst ett EP med etikett $x$ kan han bara välja att nästa EP som ska lösas är:
\begin{itemize}
  \item Det olösta EP:t med den största ID som är mindre än $x$, eller
  \item Det olösta EP:t med den minsta ID som är större än $x$.
\end{itemize}

Vid vilken tidpunkt som helst under sin resa kan Tommy äta en ego-kaka för att öka sitt ego med 1 enhet.

Vad är det minsta antalet ego-kakor Tommy måste äta för att kunna lösa alla EP:n?

\section*{Indata}
\noindent
Den första raden innehåller två heltal, $N$ $(1 \leq N \leq 5 \cdot 10^5)$ och $E$ $(0 \leq E \leq 10^9)$.

Den andra raden innehåller $N$ heltal $a_1, a_2, \dots, a_N$ $(1 \leq a_i \leq 10^9)$, där $a_i$ är svårighetsgraden
för EP:t med ID $i$.

\section*{Utdata}
\noindent
Skriv ut ett heltal: det minsta antalet ego-kakor Tommy måste äta för att kunna lösa alla EP:n.

\section*{Poängsättning}
Din lösning kommer att testas på en mängd testfallsgrupper.
För att få poäng för en grupp så måste du klara alla testfall i gruppen.

\noindent
\begin{tabular}{| l | l | p{12cm} |}
  \hline
  \textbf{Grupp} & \textbf{Poäng} & \textbf{Gränser} \\ \hline
  $1$    & $7$        & $N \leq 3$ \\ \hline
  $2$    & $12$       & $N \leq 15$ \\ \hline
  $3$    & $23$       & $N \leq 5000$ \\ \hline
  $4$    & $17$       & Det finns högst 2 olika heltal i $a_1, \dots, a_N$. \\ \hline
  $5$    & $41$       & Inga ytterligare begränsningar. \\ \hline
\end{tabular}

\section*{Förklaring av exempelfall}
Om vi löser EP:n i ordningen $a_4, a_5, a_3, a_2, a_1, a_6$ går processen till på följande sätt:  
\begin{itemize}
 \item Innan $a_4 = 2$ löses äter Tommy 1 kaka för att öka sitt ego till 2. Efter att ha löst EP:t blir hans ego 3.
 \item För $a_5 = 3$ är Tommys ego redan 3, så han löser det utan kakor. Efteråt blir hans ego 4.
 \item Innan $a_3 = 7$ löses äter Tommy 3 kakor för att höja sitt ego till 7. Efter lösning blir hans ego 8.
 \item För $a_2 = 4$ är Tommys ego 8 (inga kakor behövs). Efter lösning blir det 9.
 \item För $a_1 = 9$ löser Tommy det direkt. Hans ego ökar till 10.
 \item Slutligen, för $a_6 = 10$ löser Tommy det med sitt nuvarande ego på 10. Efteråt blir det 11.
\end{itemize}

Totalt äter Tommy 4 kakor, vilket är det minsta antal som krävs för att lösa alla EP:n.